\documentclass[../main]{subfiles}
\begin{document}
\ifSubfilesClassLoaded{\setcounter{page}{4}}{}
\section{Prefazione}
\label{sec:02_prefazione}

Nella sua critica al programma gotico, Marx scrisse: «Non ci troviamo di fronte a una società comunista che si sia sviluppata autonomamente, ma piuttosto a una società che è appena emersa dal capitalismo e che, di conseguenza, conserva in ogni sua dimensione – economica, morale e intellettuale – le tracce inconfondibili della società precedente da cui ha avuto origine». E proseguì: «Nella fase più elevata dello sviluppo comunista, una volta superata la divisione del lavoro che aliena l’individuo; quando, con essa, svanirà anche la dicotomia tra lavoro intellettuale e lavoro manuale; quando il lavoro non sarà più un mero strumento di sostentamento, ma diventerà la prima e più importante attività umana; quando, grazie allo sviluppo completo di ogni individuo, si espanderanno anche le forze produttive, e tutte le fonti della ricchezza sociale fluiranno abbondanti».Solo allora sarà possibile trascendere i limiti del diritto borghese, e la società potrà finalmente affermare: «A ciascuno secondo le sue capacità, a ciascuno secondo i suoi bisogni!»

All’ideale della proprietà privata e del dominio del capitale si contrappone l’ideale del comunismo. A una società in cui il capitale esercita il predominio, può essere contrapposta, in ultima analisi, soltanto una società comunista. Ma che cos'è, in realtà, il comunismo?

Nella presente brochure, il nostro interesse si focalizza su ciò che Marx definì la fase più elevata del comunismo: una società comunista sviluppata sulle proprie fondamenta, o il comunismo in quanto tale. La fase inferiore del comunismo, ovvero il periodo di transizione dal capitalismo al comunismo, che in seguito è stato storicamente identificato come "socialismo", sarà oggetto di analisi in seguito. Per il momento, ci concentreremo su quale debba essere, in sostanza, questa transizione, o su cosa i comunisti intendano realizzare come esito delle trasformazioni rivoluzionarie della vita sociale.

In questo contesto, ci preme analizzare le caratteristiche fondamentali che definiscono una società comunista, distinguendola da ogni forma di produzione e organizzazione sociale precedente. Lasciamo, dunque, da parte la descrizione delle tecnologie impiegate in un sistema comunista, poiché queste possono variare nel tempo (e, inevitabilmente, subiranno un’evoluzione rapida). Per fare un esempio, il capitalismo dell’era della macchina a vapore e il capitalismo dell’era delle moderne tecnologie digitali presentano differenze notevoli. Tuttavia, ciò che costituisce l’essenza del capitalismo – una società basata sulla produzione di merci, in cui la forza lavoro stessa è considerata una merce e il capitale rappresenta la relazione sociale predominante – rimane un elemento distintivo, imprescindibile e immutabile, sia nel XIX che nel XXI secolo. Comprendere questi elementi essenziali nella loro intrinseca interconnessione ci permette di delineare il concetto di capitalismo.

Analogamente, il concetto di comunismo – come quello di qualsiasi altra ideologia – comprende unicamente gli elementi imprescindibili e sufficienti a definirne le interrelazioni essenziali, affinché una società possa definirsi comunista.

La difficoltà risiede nel fatto che ci troviamo di fronte a un concetto che non ha ancora trovato piena attuazione. La storia non ha ancora conosciuto – almeno sul nostro pianeta – una società comunista sviluppata pienamente sulle proprie fondamenta (il comunismo primitivo non può essere considerato in questa categoria).

Il comunismo, tuttavia, si è ormai affermato come disciplina scientifica. L’intera evoluzione del pensiero storico ha dato i suoi frutti, generando l’idea del comunismo, e l’andamento della vita sociale, di fatto, evidenzia costantemente i limiti del capitalismo e l’ineluttabilità di una transizione verso un nuovo modello di sviluppo sociale. Le crisi di sovrapproduzione capitalistiche sono infatti accompagnate da una distruzione su scala colossale, che non si limita ai beni prodotti o alle singole imprese, ma spesso si estende anche a intere metropoli, un tempo fiorenti.

Se osserviamo la produzione sociale come un'interazione tra l'uomo e la natura, emerge con chiarezza l'irrazionalità di questo modello, che sta infliggendo al nostro pianeta danni di tale portata da renderli potenzialmente irreversibili. La produzione, e soprattutto la sovrapproduzione, di imballaggi superflui e privi di senso nel contesto di un rapporto armonioso con l'ambiente, ma imprescindibili per la commercializzazione dei prodotti, genera conseguenze devastanti. La fabbricazione di beni che non contribuiscono al benessere o al progresso dell'umanità, ma che spesso la danneggiano, e che al contempo richiedono un consumo massiccio di risorse, contaminando suolo, acqua e aria, è la cruda realtà del sistema capitalista. L'assenza di un sistema di riciclo efficiente e capillare, unito a una gestione irrazionale delle risorse naturali e delle capacità produttive, costituisce una seria minaccia di catastrofe ecologica irreversibile.

Parallelamente a questa opera di demolizione artificiale dell’ambiente che ci circonda, assistiamo a una progressiva degradazione dell’essere umano. Non si tratta unicamente della miseria più estrema in cui è relegata una parte considerevole della forza lavoro mondiale – individui costretti a sopportare la fame, a vivere in condizioni insalubri e privati dell’accesso all’istruzione e alle cure mediche – ma anche di un impoverimento relativo, nel quale vengono create nuove necessità, la cui soddisfazione si trasforma in un imperativo, una condizione imprescindibile persino per poter sopravvivere come offerenti della propria forza lavoro, pur essendo i mezzi a disposizione insufficienti per farvi fronte.Questa tendenza era già stata osservata da Engels, ma oggi, in un’epoca in cui il capitalismo non può più sopravvivere senza una produzione mirata a soddisfare i bisogni, essa si è notevolmente accentuata, assumendo proporzioni preoccupanti, soprattutto nelle aree periferiche e semi-periferiche del sistema capitalista, dove, a livello sociale, la forza lavoro non si rigenera completamente. Ciò mette seriamente a repentaglio le basi stesse della produzione capitalistica.

Ma la questione non si esaurisce qui. Per espandersi su una scala più vasta, il capitale moderno reintroduce e assoggetta a sé le forme di sfruttamento preesistenti. La schiavitù e la servitù della gleba persistono ancora oggi, poiché si rivelano convenienti e contribuiscono all’incremento del valore. Ad esempio, stando ai dati dell’OIL pubblicati nel 2014, i profitti derivanti dall’impiego di lavoro forzato a livello globale superano i 150,2 miliardi di dollari. Come si evince, il vantaggio che si trae dalla schiavitù può essere quantificato in termini strettamente capitalistici, e non solo in una forma naturale, come accadeva nell’epoca della tratta degli schiavi.

Il capitale esige incessantemente nuovi spazi per la propria espansione, e per raggiungere tale obiettivo è costretto a marginalizzare altri capitali dai mercati consolidati, riassegnando le risorse a proprio esclusivo vantaggio. Questo meccanismo genera inevitabilmente tensioni, non solo all’interno dei singoli Stati, ma soprattutto tra le diverse nazioni, sfociando in conflitti armati localizzati e aprendo la strada a una potenziale guerra di portata globale. In uno scenario in cui si manifestano segnali di una crescente instabilità (preparativi militari, dispiegamento di truppe, focolai di conflitto), e in cui esiste la minaccia incombente di armi nucleari, quali imprevedibili conseguenze potrebbero scaturire?... Inoltre, la guerra stessa paradossalmente alimenta lo sviluppo di un vasto, massiccio e spontaneo movimento anticapitalista. Resta tuttavia incerto quanto questo movimento sarà coeso e se esso riuscirà effettivamente a superare il sistema capitalistico. In definitiva, nessuno può garantire con certezza che l’umanità sopravviverà all’era del capitalismo.Ciò che emerge con chiarezza è la limitatezza di un modello di sviluppo sociale in cui il capitale esercita un predominio assoluto.

È inconfutabile che il capitalismo sia destinato a esaurirsi. La questione cruciale è se tale declino coinciderà con la fine stessa dell'umanità, oppure se l'umanità sarà in grado di adottare un nuovo paradigma esistenziale. E, soprattutto, come potrà l'umanità reinventarsi, affinché l'insieme della specie e le sue molteplici componenti possano progredire in armonia, sostenendosi reciprocamente anziché ostacolarsi? La teoria scientifica del comunismo offre una risposta a questo interrogativo, una risposta tutt'altro che superata, ma di straordinaria attualità per il futuro dell'umanità.

Ci prefiggiamo ora di delineare i punti salienti di tale teoria, una teoria sviluppata da Marx ed Engels, e successivamente arricchita dai contributi di numerosi pensatori marxisti. Tra questi, spicca in particolare l'opera teorica di Lenin.

Il significato dei termini «comunismo» e «socialismo» non si è consolidato immediatamente. Per un lungo periodo, i due termini sono stati utilizzati come sinonimi, e in alcuni casi il comunismo veniva persino considerato una fase antecedente al socialismo. Di conseguenza, è utile fare riferimento sia agli scritti di Marx ed Engels in cui compare il termine «comunismo», sia a quelli in cui viene trattato il «socialismo», a condizione che il contenuto suggerisca chiaramente che si riferiscono a una società comunista, al movimento comunista o alla teoria del comunismo. Per quanto concerne Lenin e gli altri marxisti del XX secolo, nelle loro opere, così come negli scritti successivi di Marx, la terminologia risulta già in gran parte definita.Nell’articolo intitolato “Un’iniziativa notevole”, Lenin rammenta esplicitamente ai lettori che “l’unica distinzione sostanziale tra socialismo e comunismo risiede nel fatto che il primo termine designa la fase iniziale della nuova società che si sviluppa dal capitalismo, mentre il secondo si riferisce a una fase più avanzata e successiva”. Pertanto, qualora ci avvengano di utilizzare tali termini senza fornire ulteriori chiarimenti o contestualizzazioni, essi dovranno essere interpretati in base al significato corrente.

A proposito delle reiterate argomentazioni contrarie al comunismo, costantemente avanzate da diversi apologeti ideologici della classe capitalista dominante – “il comunismo è un’utopia; abbiamo già tentato di edificare una società comunista, e l’esperimento è fallito (o non ha prodotto risultati positivi)” – replichiamo quanto segue. È indubbio che, finora, non siamo riusciti a realizzare una società comunista compiuta, ma sostenere che tale tentativo non abbia prodotto alcun beneficio è una palese falsità. Gli indiscutibili successi dell’Unione Sovietica socialista, nel suo tempo, non potevano essere negati neppure dai più accaniti anticomunisti, persino negli Stati Uniti. Chiunque conosca la storia recente del XX secolo sa bene che l’indirizzo impresso al movimento sociale verso il comunismo ha consentito a tutti i popoli che vivevano nell’URSS di raggiungere vette inimmaginabili in precedenza, durante le epoche feudale e capitalista, e tutto ciò in un arco di tempo relativamente breve.Questo è confermato sia dai dati e dalle analisi oggettive, sia dai risultati, ancora oggi insuperati, ottenuti in ambito culturale.

Se mettiamo a confronto i primi decenni di esistenza dell’organizzazione, segnati da un’evoluzione verso il comunismo, con lo stesso periodo del movimento delle “nazioni indipendenti”, sorto a seguito del crollo dell’Unione Sovietica, possiamo osservare una chiara correlazione tra l’orientamento verso il comunismo e lo sviluppo complessivo della società. Tale sviluppo non concerne soltanto, o prevalentemente, l’industria e la produzione di beni, ma anche le opportunità di crescita personale e le potenzialità umane, rese accessibili a un’ampia fetta della popolazione. \footnote{Nonostante le sue origini in un paese prevalentemente agricolo, provato dalla guerra e immerso in un contesto avverso, questo movimento ha intrapreso un cammino inesplorato, e di conseguenza la sua realizzazione è stata segnata da notevoli difficoltà, ostacoli, battute d'arresto e persino tragedie, eventi che noi comunisti non possiamo permetterci di dimenticare.Questo dinamismo sociale, permeato da marcate tendenze capitalistiche, non può essere soffocato con un semplice decreto, né può essere superato senza compromettere la produzione di beni.}

Già nel 1925, il numero degli iscritti alle scuole elementari era aumentato di ben 2.250.000 unità rispetto al 1913. Allo stesso modo, le matricole nelle scuole professionali avevano visto raddoppiarsi il loro numero. Uscita dalle devastazioni della guerra civile, la società aveva ripreso quasi immediatamente il percorso dell’istruzione. Il diritto allo studio, un tempo negato, era stato finalmente esteso a coloro che, prima del conflitto, erano stati relegati all’ignoranza. Consideriamo questo dato: l’Unione Sovietica, pur essendo giovane, conta appena un anno di esistenza, eppure ha già qualcosa di significativo da presentare al mondo capitalista, dove l’istruzione è ancora improntata a rigide distinzioni di classe. Nel sistema capitalista, infatti, un’istruzione di qualità rimaneva, e tuttora rimane, un privilegio precluso alla maggioranza della popolazione.

Nei due decenni successivi alla Rivoluzione d’Ottobre, si trovò una soluzione al problema dei bambini abbandonati: ai minori rimasti orfani a causa della Prima Guerra Mondiale e della Guerra Civile non furono offerti soltanto un riparo e del cibo, ma anche istruzione, una preparazione professionale e, di conseguenza, la prospettiva di un futuro. In netto contrasto, tra le macerie del defunto stato socialista, i bambini abbandonati sono riemersi, sebbene nella maggioranza dei casi i loro genitori siano ancora vivi. Parallelamente, il livello di istruzione generale della popolazione è calato e continua a diminuire. I sistemi educativi delle ex repubbliche sovietiche si stanno adattando alle mutate condizioni sociali, privilegiando sempre più la specializzazione e acuendo le disparità nelle opportunità formative tra i bambini provenienti da famiglie agiate e quelli provenienti da contesti più umili. Numerose testimonianze di questo processo possono essere riscontrate in Russia nel film di Konstantin Semënov ed Evgenij Spicin, intitolato «L’ultima campanella»

Nell’Unione Sovietica, il programma di alfabetizzazione – noto come «Alfabetizzazione» a partire dal 1920 – e l’istituzione di un’istruzione primaria universale («Istruzione per tutti», dal 1930) permisero di superare le difficoltà relative all’istruzione elementare in meno di due decenni, e ciò in un paese a prevalente vocazione agricola, dove, fino al 1917, oltre l’80\% della popolazione era analfabeta o scarsamente istruita. Entro il 1936, nell’URSS circa quaranta milioni di persone furono istruite, superando l’analfabetismo. Tra il 1933 e il 1937, ben oltre venti milioni di persone, tra analfabeti e individui con una preparazione limitata, frequentarono corsi presso le scuole di alfabetizzazione. Secondo il censimento del 1939, il tasso di alfabetizzazione tra la popolazione di età compresa tra i sedici e i cinquant’anni si attestava intorno al 90\%.Furono proprio quegli uomini, e non gli eroi leggendari, adornati di medaglie e mossi da una fede incrollabile (come ostentava il film «La fortezza» di Nikita Michalkov e altre opere della contemporanea produzione propagandistica), a sconfiggere la più potente, moderna e attrezzata macchina da guerra del regime nazista. All’epoca dell’invasione dell’URSS, gran parte dell’industria europea era infatti al suo servizio.

In un arco di tempo di appena vent'anni, l'analfabetismo è stato virtualmente eradicato e si è assistito alla nascita di una rete capillare di strutture sanitarie, garantendo a tutti l'accesso alle cure. Medici e insegnanti hanno iniziato a operare anche nelle aree più remote, dove prima la loro presenza era impensabile. Sono state edificate scuole, biblioteche e centri culturali, e un vigoroso impulso è stato impresso allo sviluppo scientifico, anche attraverso la realizzazione di infrastrutture adeguate. Questa opera di costruzione del comunismo, intesa come un processo volto a promuovere lo sviluppo armonioso e completo dell'individuo – attraverso l'istruzione, l'assistenza sanitaria, la scienza e la cultura – ha consentito di sconfiggere il fascismo, un'ideologia che rappresentava una minaccia per l'intero mondo, e, in seguito a una guerra cruenta e devastante, di ricostruire il paese nel minor tempo possibile.

Quali sono, invece, i fattori che hanno determinato l'abbandono delle politiche ispirate al comunismo negli anni immediatamente successivi al collasso dell'Unione Sovietica, nei territori che un tempo ne costituivano gli stati membri?

Già solo in Ucraina, dopo la dissoluzione dell'Unione Sovietica, i dati forniti dal Comitato statale per la statistica fino al 2014 indicano che la diminuzione assoluta della popolazione ha raggiunto quasi i 6,5 milioni di persone (prima del collasso, la crescita demografica era positiva). Tale dato non considera coloro che, pur essendo «ufficialmente» cittadini ucraini, sono stati costretti ad abbandonare le proprie case e a cercare lavoro all'estero. Questo fatto, di per sé, è sintomo di una catastrofe sociale di proporzioni colossali. Naturalmente, non si può dimenticare la carestia del 1932-1933. Una carestia di massa rappresenta sempre una tragedia terribile, a prescindere dalle sue reali cause. Ma anche secondo le stime più prudenti (e sottolineo stime, e non mere affermazioni ideologiche) elaborate dall'Istituto di demografia dell'Accademia nazionale delle scienze dell'Ucraina, la carestia in Ucraina ha causato la morte non di 7 o 10 milioni, bensì di 3,9 milioni di persone.Non si è tenuto conto unicamente delle perdite demografiche, ma si è effettuato un confronto tra il numero effettivo di abitanti e le stime previsionali, considerando il notevole incremento registrato negli anni immediatamente precedenti e successivi alla carestia. Qualora si calcolassero le perdite assolute della popolazione ucraina in seguito al collasso dell'Unione Sovietica, applicando una metodologia analoga, il quadro risulterebbe ben più preoccupante rispetto a quello che si ottiene sottraendo i quasi 6,5 milioni di persone.

Analogamente, anche in Lituania, Lettonia ed Estonia si osserva un sensibile declino demografico. In Russia, grazie ai flussi migratori, il calo demografico è stato meno accentuato, passando da 148,6 milioni nel 1991 a 144,3 milioni nel 2016.

Se si prendesse in considerazione la direzione dei flussi migratori e la composizione del mercato del lavoro nelle stesse repubbliche dell'Asia centrale, si evincerebbe chiaramente che la qualità della vita complessiva, unitamente alle opportunità per i cittadini di tali repubbliche di trovare un impiego che consenta loro di mantenere la propria famiglia nel paese d'origine, è notevolmente peggiorata.

I conflitti che insanguinano il territorio dell’ex Unione Sovietica sono purtroppo divenuti una costante: dalla contesa tra Azerbaigian e Armenia per il Nagorno-Karabakh – un conflitto nato in epoca sovietica, ma che ha trovato la sua più cruenta espressione nel periodo immediatamente successivo al collasso dell’URSS – alla guerra civile in Tagikistan, passando per la questione della Transnistria, le due guerre cecene, i conflitti in Georgia e Abkhazia, in Ossezia del Sud e, infine, nel Donbass. In tutte queste guerre, gli esseri umani si sono massacrati a vicenda, e in ultima analisi, a trarne vantaggio sono stati i gruppi di interesse e le élite dominanti, i veri sfruttatori del popolo.

Si potrebbe obiettare che anche la guerra civile del 1917 fu cruenta e, secondo diverse stime, causò tra gli otto e i tredici milioni di morti da entrambe le parti. Tale cifra include non solo i caduti in battaglia, ma anche coloro che perirono a causa di malattie, fame e terrore. Certamente, è un fatto. E, naturalmente, sarebbe stato preferibile evitare il conflitto. Tuttavia, in primo luogo, la guerra non fu scatenata dai bolscevichi né dal proletariato, bensì dai rappresentanti delle vecchie classi dominanti, riluttanti a rinunciare al loro potere. In secondo luogo, in tale guerra i lavoratori difendevano i propri interessi, e non quelli di coloro che si arricchiscono sfruttando il loro lavoro. Infine, tale conflitto rappresentò una lotta per la possibilità di realizzare una trasformazione radicale dell'intera società, e non una semplice ridistribuzione delle aree di influenza del capitale.È bene trattare a parte la Grande Guerra Patriottica, poiché essa rappresentò uno scontro decisivo contro il fascismo, nel corso del quale il popolo sovietico difese non solo i propri interessi, ma anche quelli dell’intera umanità.

Come si può facilmente constatare, un semplice confronto tra i primi decenni di esistenza dell’URSS e le condizioni di vita nelle repubbliche sorte dopo la sua dissoluzione offre spunti interessanti per riflettere sui vantaggi di una società in cui sia presente un’aspirazione, una spinta verso il comunismo. Ciò, nonostante il percorso non sia stato lineare e numerosi aspetti non possano in alcun modo essere definiti comunisti. Si tratta, piuttosto, di processi e fenomeni intrinsecamente capitalisti sviluppatisi all’interno di una società socialista. Sono proprio questi elementi che i detrattori dell’URSS criticano aspramente, sebbene, quando tale critica è pienamente giustificata, essa si riveli in realtà un’autocritica nei confronti del mondo capitalista.

A soli dodici anni dalla Grande Guerra Patriottica, un conflitto che aveva mietuto milioni di vite e devastato ampie porzioni del territorio, l’Unione Sovietica, spinta dalla sua ideologia comunista, riuscì a lanciare in orbita il primo satellite artificiale. Poco tempo dopo, un uomo varcò la soglia dello spazio. Questo straordinario progresso scientifico fu reso possibile dalla visione comunista, un modello che, nel contesto attuale, sembra più proiettato verso il futuro che ancorato al passato, poiché la sua piena attuazione si scontra inevitabilmente con il principio della proprietà privata dei mezzi di produzione. L’elenco dei successi potrebbe estendersi ulteriormente, includendo non solo le conquiste dell’Unione Sovietica, ma anche quelle di altri paesi che si sono ispirati al modello socialista delineato da Lenin, immaginando un futuro in cui la rivoluzione mondiale avrebbe trionfato.

Naturalmente, in questa grande impresa edilizia ci sono stati non solo successi, ma anche battute d’arresto e fallimenti. Ma, anche a prescindere da quanto detto finora, l’argomento

“non siamo riusciti a costruire il comunismo, quindi è impossibile e niente di tutto ciò avrà successo” presenta lo stesso difetto logico che affliggeva le affermazioni degli oppositori delle idee alla base, ad esempio, dello sviluppo e dell’evoluzione dell’aviazione. Questi ultimi sostenevano che era impossibile costruire aeromobili più pesanti dell’aria. È come affermare che un individuo o l’umanità nel suo complesso non saranno in grado di imparare qualcosa solo perché finora non ci sono riusciti, nonostante i tentativi.

Argomentazioni tanto illogiche denotano una superficialità nell’analisi della questione e delle relative obiezioni. Esse non intaccano, pertanto, minimamente le fondamenta del pensiero comunista. Il fallimento del primo tentativo storico di realizzare il comunismo non fa altro che sottolineare l’urgenza di un’indagine approfondita sull’esperienza maturata, al fine di comprenderne le cause e di vagliare con maggiore attenzione le condizioni necessarie per una proficua transizione verso tale sistema. L’esistenza stessa e la struttura dell’ideale comunista, nato come riflesso delle dinamiche e delle contraddizioni insite nella società capitalista, lo esigono.

Dunque, poiché il comunismo, tra le altre cose, si è affermato come disciplina scientifica, iniziamo a studiarlo. La presente brochure, ovviamente, non ambisce a essere un'analisi esaustiva e approfondita del comunismo. Piuttosto, essa si propone di illustrare i principi cardine della scienza comunista, così come sono stati elaborati nella teoria marxista. Un'indagine più ampia e dettagliata su tale tema rappresenta un compito per il futuro prossimo. Invitiamo pertanto i nostri lettori a partecipare attivamente a questo progetto, anche attraverso la critica costruttiva del nostro lavoro.
\end{document}
